\begin{frame}[standout]

    Avez-vous été attentifs ?

    À vous de jouer !

\end{frame}

{
\setbeamercovered{transparent}

\begin{frame}{Q1}

    \textbf{Quelle affirmation est vraie à propos du modèle relationnel ?}

    \begin{itemize}
        \item \uncover<1>{\textcolor{kh-red}{Utilise des graphes pour représenter les relations entre les données.}}
        \item \uncover<1>{\textcolor{kh-blue}{Stocke les données en documents JSON imbriqués.}}
        \item \textbf<2>{\textcolor{kh-green}{Organise les données en tables interconnectées par des clés.}}
        \item \uncover<1>{\textcolor{kh-yellow}{Se base sur des arbres hiérarchiques pour organiser les données.}}
    \end{itemize}

\end{frame}

\begin{frame}{Q2}

    \textbf{Quel est le principal langage utilisé pour interagir avec une base de données relationnelle ?}

    \begin{itemize}
        \item \uncover<1>{\textcolor{kh-red}{HTML}}
        \item \textbf<2>{\textcolor{kh-blue}{SQL}}
        \item \uncover<1>{\textcolor{kh-green}{JavaScript}}
        \item \uncover<1>{\textcolor{kh-yellow}{CSS}}
    \end{itemize}

\end{frame}

\begin{frame}{Q3}

    \textbf{Dans une base de données relationnelle, que représente une clé primaire ?}

    \begin{itemize}
        \item \uncover<1>{\textcolor{kh-red}{Un champ qui est utilisé pour indexer les données}}
        \item \uncover<1>{\textcolor{kh-blue}{Une colonne qui permet de lier plusieurs tables entre elles}}
        \item \uncover<1>{\textcolor{kh-green}{Une méthode pour effectuer des opérations de sauvegarde}}
        \item \textbf<2>{\textcolor{kh-yellow}{Une contrainte garantissant que chaque enregistrement est unique}}
    \end{itemize}

\end{frame}

\begin{frame}{Q4}

    \textbf{Que veut dire SQL ?}

    \begin{itemize}
        \item \uncover<1>{\textcolor{kh-red}{Systematic Query Layer}}
        \item \textbf<2>{\textcolor{kh-blue}{Structured Query Language}}
        \item \uncover<1>{\textcolor{kh-green}{Standard Query Language}}
        \item \uncover<1>{\textcolor{kh-yellow}{Simple Query Log}}
    \end{itemize}

\end{frame}

\begin{frame}{Q5}

    \textbf{Quel est l'objectif principal de la \textbf{normalisation} dans les bases de données relationnelles ?}

    \begin{itemize}
        \item \textbf<2>{\textcolor{kh-red}{Réduire la redondance et améliorer l'intégrité des données}}
        \item \uncover<1>{\textcolor{kh-blue}{Accélérer les requêtes}}
        \item \uncover<1>{\textcolor{kh-green}{Simplifier la syntaxe SQL}}
        \item \uncover<1>{\textcolor{kh-yellow}{Assurer la compatibilité avec différents systèmes de gestion de bases de données}}
    \end{itemize}

\end{frame}

\begin{frame}{Q6}

    \textbf{Qu'est-ce qu'une jointure dans le contexte des bases de données relationnelles ?}

    \begin{itemize}
        \item \uncover<1>{\textcolor{kh-red}{Opération pour fusionner deux bases de données distinctes}}
        \item \uncover<1>{\textcolor{kh-blue}{Procédure pour créer des tables dans une base de données}}
        \item \uncover<1>{\textcolor{kh-green}{Mécanisme pour sauvegarder des données dans une base de données}}
        \item \textbf<2>{\textcolor{kh-yellow}{Requête combinant les lignes de plusieurs tables selon une condition}}
    \end{itemize}

\end{frame}

\begin{frame}{Q7}

    \textbf{Dans un système de base de bases de données relationnelles, une relation entre deux tables est souvent établie par l'utilisation d'une clé étrangère, qui fait référence à une clé primaire d'une autre table. Cependant une clé étrangère ne peut pas contenir de valeurs nulles.}

    \begin{itemize}
        \item Vrai
        \item \textcolor<2>{kh-green}{\textbf<2>{Faux}}
    \end{itemize}

\end{frame}

}
